\documentclass[11pt,a4paper]{article}
\usepackage[T1]{fontenc}
\usepackage[utf8]{inputenc}
\usepackage[brazil]{babel}
\usepackage{amsmath,amssymb,textcomp}
\usepackage[margin=2.2cm]{geometry}
\usepackage{enumitem}
\usepackage{array}
\setlength{\parindent}{0pt}
\setlength{\parskip}{0.6em}

\begin{document}
\section*{Avaliação de questões --- avaliador 2}
\subsection*{Como responder}

São 22 questões de Matemática do Ensino Médio. Leia cada uma como leria uma
questão que você pensasse em usar com a sua turma, e responda às quatro
perguntas que vêm logo abaixo dela.

Onde se pergunta pelo \textbf{nível de dificuldade}, responda pensando nos
seus alunos, não num aluno ideal. É essa resposta que permite calibrar o
sistema: ele classifica cada questão como fácil, média ou difícil por conta
própria, e precisamos saber o quanto isso corresponde à sala de aula real.

\textbf{A amostra tem qualidade variável, e pode conter questões com erro
matemático.} Se encontrar algum, aponte onde.

Não há resposta certa sobre você: o que está sendo avaliado é o sistema que
produziu estas questões, não quem as lê. Um ``recusada'' bem justificado vale
mais para este trabalho do que um ``aceita'' por gentileza.

Você pode responder direto neste documento impresso ou na planilha
\texttt{respostas-avaliador-2.csv}, que já vem com os códigos dos itens nas linhas.

Tempo estimado: cerca de 75 minutos.
\bigskip\hrule\bigskip
\begin{samepage}
\subsection*{Q035}
Ao girar um ponto sobre o ciclo trigonométrico (circunferência de raio 1 centrada na origem), sua ordenada varia continuamente entre $-1$ e $1$, repetindo os mesmos valores a cada volta completa de $2\pi$ radianos. Essa ordenada é justamente o valor de $\sin(x)$ representado no plano cartesiano. Considere agora a função $f(x) = 3\sin(x) - 1$, definida para todo número real $x$. Com base na relação entre o que ocorre no ciclo trigonométrico e o comportamento do gráfico de $f$ no plano cartesiano, assinale a alternativa que indica corretamente o domínio, o período e a imagem de $f$.
\begin{enumerate}[label=(\alph*),nosep,leftmargin=*]
\item Domínio $\mathbb{R}$, período $2\pi$, imagem $[-4, 2]$
\item Domínio $\mathbb{R}$, período $2\pi$, imagem $[-1, 1]$
\item Domínio $\mathbb{R}$, período $6\pi$, imagem $[-4, 2]$
\item Domínio $\mathbb{R}$, período $2\pi$, imagem $[-2, 4]$
\end{enumerate}

\textbf{Gabarito proposto:} Domínio $\mathbb{R}$, período $2\pi$, imagem $[-4, 2]$

\textbf{Habilidade declarada:} EM13MAT404 --- Identificar as características fundamentais das funções seno e cosseno (periodicidade, domínio, imagem), por meio da comparação das representações em ciclos trigonométricos e em planos cartesianos, com ou sem apoio de tecnologias digitais.
\end{samepage}

\noindent\rule{\linewidth}{0.4pt}
\textbf{Tem erro matemático?}\quad
\mbox{$\square$~não} \quad \mbox{$\square$~sim, aqui:} \underline{\hspace{5cm}}
\quad \mbox{$\square$~não sei dizer}

\textbf{Que nível de dificuldade esta questão tem para os seus alunos?}\quad
\mbox{$\square$~fácil} \quad \mbox{$\square$~média} \quad
\mbox{$\square$~difícil} \quad \mbox{$\square$~fora do alcance da turma}

\textbf{Você usaria esta questão?}\quad
\mbox{$\square$~aceita} \quad \mbox{$\square$~aceita com ajuste} \quad
\mbox{$\square$~recusada}

\textbf{Por quê?} \emph{(obrigatório quando não for ``aceita'')}

\underline{\hspace{\linewidth}}

\underline{\hspace{\linewidth}}
\vspace{0.6em}

\begin{samepage}
\subsection*{Q009}
Uma locadora de bicicletas cobra pelo aluguel um valor que combina uma taxa fixa de retirada com um valor proporcional ao número de horas de uso. Um cliente que utilizou a bicicleta por 2 horas pagou R\$\,18,00, enquanto outro cliente, que a utilizou por 5 horas, pagou R\$\,33,00.

Buscando atrair clientes que alugam por períodos mais longos, a locadora lançou um novo plano de cobrança, mantendo a mesma taxa fixa de retirada, mas alterando o valor cobrado por hora de uso. Nesse novo plano, um cliente que utiliza a bicicleta por 8 horas paga R\$\,10,00 a menos do que pagaria, nesse mesmo período de uso, pelo plano antigo.

Determine a lei que descreve o custo do novo plano em função do tempo de uso $t$ (em horas) e avalie para quais valores de $t$ esse novo plano é financeiramente mais vantajoso do que o plano antigo.
\begin{enumerate}[label=(\alph*),nosep,leftmargin=*]
\item D(t) = 3,75t + 8; o novo plano é mais vantajoso, com custo estritamente menor que o antigo, para todo tempo de uso t > 0 (os dois planos coincidem apenas em t = 0).
\item D(t) = 5t - 2; o novo plano é mais vantajoso para todo t > 0.
\item D(t) = 4,75t + 8; o novo plano é mais vantajoso para todo t > 0.
\item D(t) = 3,75t + 8; porém o novo plano só se torna mais vantajoso para t > 8 horas.
\end{enumerate}

\textbf{Gabarito proposto:} D(t) = 3,75t + 8; o novo plano é mais vantajoso (custo estritamente menor) para todo t > 0, sendo os planos equivalentes apenas em t = 0.

\textbf{Habilidade declarada:} EM13MAT302 --- Resolver e elaborar problemas cujos modelos são as funções polinomiais de 1º e 2º graus, em contextos diversos, incluindo ou não tecnologias digitais.
\end{samepage}

\noindent\rule{\linewidth}{0.4pt}
\textbf{Tem erro matemático?}\quad
\mbox{$\square$~não} \quad \mbox{$\square$~sim, aqui:} \underline{\hspace{5cm}}
\quad \mbox{$\square$~não sei dizer}

\textbf{Que nível de dificuldade esta questão tem para os seus alunos?}\quad
\mbox{$\square$~fácil} \quad \mbox{$\square$~média} \quad
\mbox{$\square$~difícil} \quad \mbox{$\square$~fora do alcance da turma}

\textbf{Você usaria esta questão?}\quad
\mbox{$\square$~aceita} \quad \mbox{$\square$~aceita com ajuste} \quad
\mbox{$\square$~recusada}

\textbf{Por quê?} \emph{(obrigatório quando não for ``aceita'')}

\underline{\hspace{\linewidth}}

\underline{\hspace{\linewidth}}
\vspace{0.6em}

\begin{samepage}
\subsection*{Q008}
Uma companhia de saneamento calcula o valor mensal da conta de água residencial, em reais, de acordo com o consumo mensal $x$ (em metros cúbicos), conforme a tabela abaixo:

\begin{center}
\begin{tabular}{|p{0.430\linewidth}|p{0.430\linewidth}|}
\hline
\textbf{Faixa de consumo} & \textbf{Valor cobrado} \\
\hline
$0 \le x \le 10$ & Taxa mínima fixa de R\$\,25,00 \\
\hline
$10 < x \le 25$ & R\$\,25,00 mais R\$\,2,50 por cada $m^3$ que exceder 10 $m^3$ \\
\hline
$x > 25$ & R\$\,62,50 mais R\$\,4,00 por cada $m^3$ que exceder 25 $m^3$ \\
\hline
\end{tabular}
\end{center}

Seja $V(x)$ o valor da conta, em reais, para um consumo de $x$ $m^3$, com $x \ge 0$. Qual das alternativas descreve corretamente o gráfico de $V(x)$?
\begin{enumerate}[label=(\alph*),nosep,leftmargin=*]
\item Para $0\le x\le10$, o gráfico é um segmento horizontal na altura $25$. Para $10<x\le25$, um segmento crescente que parte do ponto $(10;25)$ e chega ao ponto $(25;62{,}5)$, com inclinação $2{,}5$. Para $x>25$, uma semirreta crescente que parte de $(25;62{,}5)$, com inclinação $4$, sem limite superior.
\item Para $0\le x\le10$, segmento horizontal em $25$. Para $10<x\le25$, segmento crescente que parte do ponto $(10;0)$ e chega a $(25;37{,}5)$, com inclinação $2{,}5$. Para $x>25$, semirreta que parte de $(25;37{,}5)$, com inclinação $4$.
\item Para $0\le x\le10$, segmento horizontal em $25$. Para $10<x\le25$, segmento crescente que parte de $(10;25)$ e chega a $(25;85)$, com inclinação $4$. Para $x>25$, semirreta que parte de $(25;85)$, com inclinação $2{,}5$.
\item Para $0\le x\le10$, o gráfico é um segmento crescente que parte da origem $(0;0)$ até $(10;25)$, com inclinação $2{,}5$; a partir daí, o gráfico segue como nas demais faixas descritas na tabela.
\end{enumerate}

\textbf{Gabarito proposto:} A

\textbf{Habilidade declarada:} EM13MAT405 --- Reconhecer funções definidas por uma ou mais sentenças (como a tabela do Imposto de Renda, contas de luz, água, gás etc.), em suas representações algébrica e gráfica, convertendo essas representações de uma para outra e identificando domínios de validade, imagem, crescimento e decrescimento.
\end{samepage}

\noindent\rule{\linewidth}{0.4pt}
\textbf{Tem erro matemático?}\quad
\mbox{$\square$~não} \quad \mbox{$\square$~sim, aqui:} \underline{\hspace{5cm}}
\quad \mbox{$\square$~não sei dizer}

\textbf{Que nível de dificuldade esta questão tem para os seus alunos?}\quad
\mbox{$\square$~fácil} \quad \mbox{$\square$~média} \quad
\mbox{$\square$~difícil} \quad \mbox{$\square$~fora do alcance da turma}

\textbf{Você usaria esta questão?}\quad
\mbox{$\square$~aceita} \quad \mbox{$\square$~aceita com ajuste} \quad
\mbox{$\square$~recusada}

\textbf{Por quê?} \emph{(obrigatório quando não for ``aceita'')}

\underline{\hspace{\linewidth}}

\underline{\hspace{\linewidth}}
\vspace{0.6em}

\begin{samepage}
\subsection*{Q032}
Uma roda-gigante de um parque de diversões tem raio de 15 m, e seu eixo central está fixado a 17 m de altura em relação ao solo. A roda gira com velocidade angular constante, completando uma volta a cada 40 segundos. No instante t = 0 (em segundos, contado a partir da abertura do parque), uma cabine específica está exatamente no ponto mais baixo de sua trajetória circular. O parque mantém essa roda-gigante em funcionamento ininterrupto durante as 12 horas de seu horário de atendimento diário (das 9h às 21h).

Um estudante decide modelar a altura h(t), em metros, dessa cabine em relação ao solo, t segundos após a abertura do parque, representando simultaneamente o movimento da cabine no ciclo trigonométrico (posição angular ao longo da circunferência) e no plano cartesiano (gráfico de h em função de t).

a) Determine a lei da função h(t), explicitando como o ângulo percorrido no ciclo trigonométrico se relaciona com a altura registrada no gráfico cartesiano.

b) Indique o período dessa função e explique seu significado físico no contexto da roda-gigante.

c) Considerando que a roda-gigante funciona apenas durante o horário diário de operação do parque, determine o domínio de h(t) para um dia de funcionamento.

d) Determine a imagem de h(t), justificando o resultado a partir da comparação entre a variação do cosseno no ciclo trigonométrico (entre -1 e 1) e a amplitude observada no gráfico cartesiano.

\textbf{Gabarito proposto:} h(t) = 17 - 15·cos($\pi$t/20) (equivalente a 17 + 15·sen($\pi$t/20 - $\pi$/2)); período = 40 s; domínio, para um dia de operação = [0, 43200] segundos; imagem = [2, 32] metros.

\textbf{Habilidade declarada:} EM13MAT404 --- Identificar as características fundamentais das funções seno e cosseno (periodicidade, domínio, imagem), por meio da comparação das representações em ciclos trigonométricos e em planos cartesianos, com ou sem apoio de tecnologias digitais.
\end{samepage}

\noindent\rule{\linewidth}{0.4pt}
\textbf{Tem erro matemático?}\quad
\mbox{$\square$~não} \quad \mbox{$\square$~sim, aqui:} \underline{\hspace{5cm}}
\quad \mbox{$\square$~não sei dizer}

\textbf{Que nível de dificuldade esta questão tem para os seus alunos?}\quad
\mbox{$\square$~fácil} \quad \mbox{$\square$~média} \quad
\mbox{$\square$~difícil} \quad \mbox{$\square$~fora do alcance da turma}

\textbf{Você usaria esta questão?}\quad
\mbox{$\square$~aceita} \quad \mbox{$\square$~aceita com ajuste} \quad
\mbox{$\square$~recusada}

\textbf{Por quê?} \emph{(obrigatório quando não for ``aceita'')}

\underline{\hspace{\linewidth}}

\underline{\hspace{\linewidth}}
\vspace{0.6em}

\begin{samepage}
\subsection*{Q021}
No mesmo plano cartesiano estão representadas duas curvas, $C_1$ e $C_2$, associadas a duas funções reais $f$ e $g$, cada uma definida por uma única sentença (sem sentenças por partes). Sabe-se que:

\begin{itemize}[nosep,leftmargin=*]
\item $C_1$ passa pelos pontos $\left(-1,\ \dfrac{1}{2}\right)$, $(0,1)$, $(1,2)$ e $(2,4)$;
\item $C_2$ passa pelos pontos $\left(\dfrac{1}{2},\ -1\right)$, $(1,0)$, $(2,1)$ e $(4,2)$;
\item $C_1$ e $C_2$ são simétricas uma em relação à outra em torno da reta $y=x$.
\end{itemize}

Com base na leitura desses pontos e na relação de simetria entre as curvas, assinale a alternativa que descreve corretamente o domínio, a imagem e o comportamento de crescimento de $f$ e de $g$.
\begin{enumerate}[label=(\alph*),nosep,leftmargin=*]
\item $f$ tem domínio $\mathbb{R}$ e imagem $(0,+\infty)$; $g$ tem domínio $(0,+\infty)$ e imagem $\mathbb{R}$; ambas são crescentes em todo o seu domínio.
\item $f$ tem domínio $(0,+\infty)$ e imagem $\mathbb{R}$; $g$ tem domínio $\mathbb{R}$ e imagem $(0,+\infty)$; ambas são crescentes.
\item $f$ tem domínio $\mathbb{R}$ e imagem $(0,+\infty)$; $g$ tem domínio $\mathbb{R}$ e imagem $\mathbb{R}$; ambas são crescentes.
\item $f$ é crescente, mas $g$ é decrescente, pois $C_2$ passa por pontos com ordenada negativa.
\end{enumerate}

\textbf{Gabarito proposto:} Alternativa A: $f$ tem domínio $\mathbb{R}$ e imagem $(0,+\infty)$; $g$ tem domínio $(0,+\infty)$ e imagem $\mathbb{R}$; ambas são crescentes em todo o seu domínio.

\textbf{Habilidade declarada:} EM13MAT403 --- Comparar e analisar as representações, em plano cartesiano, das funções exponencial e logarítmica para identificar as características fundamentais (domínio, imagem, crescimento) de cada uma, com ou sem apoio de tecnologias digitais, estabelecendo relações entre elas.
\end{samepage}

\noindent\rule{\linewidth}{0.4pt}
\textbf{Tem erro matemático?}\quad
\mbox{$\square$~não} \quad \mbox{$\square$~sim, aqui:} \underline{\hspace{5cm}}
\quad \mbox{$\square$~não sei dizer}

\textbf{Que nível de dificuldade esta questão tem para os seus alunos?}\quad
\mbox{$\square$~fácil} \quad \mbox{$\square$~média} \quad
\mbox{$\square$~difícil} \quad \mbox{$\square$~fora do alcance da turma}

\textbf{Você usaria esta questão?}\quad
\mbox{$\square$~aceita} \quad \mbox{$\square$~aceita com ajuste} \quad
\mbox{$\square$~recusada}

\textbf{Por quê?} \emph{(obrigatório quando não for ``aceita'')}

\underline{\hspace{\linewidth}}

\underline{\hspace{\linewidth}}
\vspace{0.6em}

\begin{samepage}
\subsection*{Q065}
Um oceanógrafo registra a altura da maré (em metros) em um pequeno porto ao longo de um dia. Ele observa que a maré varia de forma periódica e simétrica em torno de um nível médio, atingindo seu valor máximo de 1,8 m exatamente à meia-noite ($t=0$, com $t$ em horas) e seu valor mínimo de 0,2 m exatamente às 6h da manhã. Esse padrão se repete a cada 12 horas.

a) Modele a altura da maré, em metros, por uma função do tipo $h(t) = A\cos(Bt) + D$, com $t$ em horas contado a partir da meia-noite. Determine explicitamente os valores de $A$, $B$ e $D$, justificando cada um a partir dos dados do fenômeno (amplitude, período e deslocamento vertical do gráfico em relação ao eixo $t$).

b) Determine todos os horários dentro das primeiras 12 horas do dia (isto é, para $0 \le t < 12$) em que a maré atinge exatamente 1,4 m de altura. Para cada horário encontrado, indique se a maré está subindo ou descendo naquele instante, justificando sua resposta apenas a partir do comportamento gráfico da função cosseno (sem usar derivadas).

\textbf{Gabarito proposto:} h(t) = 0,8·cos($\pi$t/6) + 1,0 (A=0,8 m; B=$\pi$/6 rad/h; D=1,0 m). A maré atinge 1,4 m às 2h (descendo) e às 10h (subindo).

\textbf{Habilidade declarada:} EM13MAT306 --- Resolver e elaborar problemas em contextos que envolvem fenômenos periódicos reais, como ondas sonoras, ciclos menstruais, movimentos cíclicos, entre outros, e comparar suas representações com as funções seno e cosseno, no plano cartesiano, com ou sem apoio de aplicativos de álgebra e geometria.
\end{samepage}

\noindent\rule{\linewidth}{0.4pt}
\textbf{Tem erro matemático?}\quad
\mbox{$\square$~não} \quad \mbox{$\square$~sim, aqui:} \underline{\hspace{5cm}}
\quad \mbox{$\square$~não sei dizer}

\textbf{Que nível de dificuldade esta questão tem para os seus alunos?}\quad
\mbox{$\square$~fácil} \quad \mbox{$\square$~média} \quad
\mbox{$\square$~difícil} \quad \mbox{$\square$~fora do alcance da turma}

\textbf{Você usaria esta questão?}\quad
\mbox{$\square$~aceita} \quad \mbox{$\square$~aceita com ajuste} \quad
\mbox{$\square$~recusada}

\textbf{Por quê?} \emph{(obrigatório quando não for ``aceita'')}

\underline{\hspace{\linewidth}}

\underline{\hspace{\linewidth}}
\vspace{0.6em}

\begin{samepage}
\subsection*{Q087}
Considere a progressão aritmética $(a_n)$ com primeiro termo $a_1 = 7$ e razão $r = 4$: $(7, 11, 15, 19, 23, \dots)$. Essa PA pode ser vista como a restrição da função afim $f(x) = 4x + 3$ ao conjunto dos números naturais não nulos, de modo que $a_n = f(n)$ para todo $n \geq 1$ (o gráfico da PA é o conjunto de pontos discretos $(n, a_n)$ sobre a reta que representa $f$).

Qual é o menor termo dessa PA que é maior do que 100?
\begin{enumerate}[label=(\alph*),nosep,leftmargin=*]
\item 103
\item 99
\item 107
\item 100
\end{enumerate}

\textbf{Gabarito proposto:} 103

\textbf{Habilidade declarada:} EM13MAT507 --- Identificar e associar sequências numéricas (PA) a funções afins de domínios discretos para análise de propriedades, incluindo dedução de algumas fórmulas e resolução de problemas.
\end{samepage}

\noindent\rule{\linewidth}{0.4pt}
\textbf{Tem erro matemático?}\quad
\mbox{$\square$~não} \quad \mbox{$\square$~sim, aqui:} \underline{\hspace{5cm}}
\quad \mbox{$\square$~não sei dizer}

\textbf{Que nível de dificuldade esta questão tem para os seus alunos?}\quad
\mbox{$\square$~fácil} \quad \mbox{$\square$~média} \quad
\mbox{$\square$~difícil} \quad \mbox{$\square$~fora do alcance da turma}

\textbf{Você usaria esta questão?}\quad
\mbox{$\square$~aceita} \quad \mbox{$\square$~aceita com ajuste} \quad
\mbox{$\square$~recusada}

\textbf{Por quê?} \emph{(obrigatório quando não for ``aceita'')}

\underline{\hspace{\linewidth}}

\underline{\hspace{\linewidth}}
\vspace{0.6em}

\begin{samepage}
\subsection*{Q024}
A tabela abaixo relaciona valores de $x$ e os correspondentes valores de $y$, obtidos a partir de uma mesma regra de formação (os valores de $x$ não variam em intervalos iguais):

\begin{center}
\begin{tabular}{|p{0.172\linewidth}|p{0.172\linewidth}|p{0.172\linewidth}|p{0.172\linewidth}|p{0.172\linewidth}|}
\hline
\textbf{$x$} & \textbf{$-3$} & \textbf{$1$} & \textbf{$4$} & \textbf{$8$} \\
\hline
$y$ & $14$ & $2$ & $-7$ & $-19$ \\
\hline
\end{tabular}
\end{center}

Um estudante, analisando os pares ordenados $(x,y)$ da tabela, conjectura que essa relação pode ser representada por uma função polinomial do 1º grau. Assinale a alternativa que apresenta a lei de formação $f(x)$ que generaliza corretamente o padrão observado na tabela.
\begin{enumerate}[label=(\alph*),nosep,leftmargin=*]
\item $f(x) = -3x + 5$
\item $f(x) = -12x + 14$
\item $f(x) = 3x - 1$
\item $f(x) = -11x + 13$
\end{enumerate}

\textbf{Gabarito proposto:} f(x) = -3x + 5

\textbf{Habilidade declarada:} EM13MAT501 --- Investigar relações entre números expressos em tabelas para representá-los no plano cartesiano, identificando padrões e criando conjecturas para generalizar e expressar algebricamente essa generalização, reconhecendo quando essa representação é de função polinomial de 1º grau.
\end{samepage}

\noindent\rule{\linewidth}{0.4pt}
\textbf{Tem erro matemático?}\quad
\mbox{$\square$~não} \quad \mbox{$\square$~sim, aqui:} \underline{\hspace{5cm}}
\quad \mbox{$\square$~não sei dizer}

\textbf{Que nível de dificuldade esta questão tem para os seus alunos?}\quad
\mbox{$\square$~fácil} \quad \mbox{$\square$~média} \quad
\mbox{$\square$~difícil} \quad \mbox{$\square$~fora do alcance da turma}

\textbf{Você usaria esta questão?}\quad
\mbox{$\square$~aceita} \quad \mbox{$\square$~aceita com ajuste} \quad
\mbox{$\square$~recusada}

\textbf{Por quê?} \emph{(obrigatório quando não for ``aceita'')}

\underline{\hspace{\linewidth}}

\underline{\hspace{\linewidth}}
\vspace{0.6em}

\begin{samepage}
\subsection*{Q057}
A magnitude $M$ de um terremoto na escala Richter relaciona-se com a energia $E$ (em joules) liberada pelo abalo por meio da expressão $M = \dfrac{2}{3}\log_{10}\left(\dfrac{E}{E_0}\right)$, em que $E_0 = 10^{4,4}$ J é uma constante de referência. Uma cidade foi atingida por um terremoto de magnitude 5,0. Meses depois, a mesma região sofreu outro abalo, de magnitude 6,5. Aproximadamente quantas vezes maior foi a energia liberada pelo segundo terremoto em comparação com o primeiro?
\begin{enumerate}[label=(\alph*),nosep,leftmargin=*]
\item Aproximadamente 178 vezes maior.
\item Aproximadamente 1,3 vezes maior (30\% a mais), pois $6{,}5 \div 5{,}0 = 1{,}3$.
\item Aproximadamente 32 vezes maior, pois $10^{1,5} \approx 31{,}6$.
\item Aproximadamente 10 vezes maior, pois $10^{1} = 10$.
\end{enumerate}

\textbf{Gabarito proposto:} Aproximadamente 178 vezes (10\textasciicircum{}2,25)

\textbf{Habilidade declarada:} EM13MAT305 --- Resolver e elaborar problemas com funções logarítmicas nos quais é necessário compreender e interpretar a variação das grandezas envolvidas, em contextos como os de abalos sísmicos, pH, radioatividade, Matemática Financeira, entre outros.
\end{samepage}

\noindent\rule{\linewidth}{0.4pt}
\textbf{Tem erro matemático?}\quad
\mbox{$\square$~não} \quad \mbox{$\square$~sim, aqui:} \underline{\hspace{5cm}}
\quad \mbox{$\square$~não sei dizer}

\textbf{Que nível de dificuldade esta questão tem para os seus alunos?}\quad
\mbox{$\square$~fácil} \quad \mbox{$\square$~média} \quad
\mbox{$\square$~difícil} \quad \mbox{$\square$~fora do alcance da turma}

\textbf{Você usaria esta questão?}\quad
\mbox{$\square$~aceita} \quad \mbox{$\square$~aceita com ajuste} \quad
\mbox{$\square$~recusada}

\textbf{Por quê?} \emph{(obrigatório quando não for ``aceita'')}

\underline{\hspace{\linewidth}}

\underline{\hspace{\linewidth}}
\vspace{0.6em}

\begin{samepage}
\subsection*{Q034}
Duas empresas de aplicativo de transporte, A e B, calculam o valor da corrida em função da distância percorrida $x$ (em km). A empresa A cobra R\$\,2,50 por quilômetro rodado, sem nenhuma taxa adicional. A empresa B cobra uma taxa fixa de embarque de R\$\,3,00 mais R\$\,2,00 por quilômetro rodado. Sejam $y_A(x)$ e $y_B(x)$, em reais, os custos totais das corridas das empresas A e B, respectivamente, em função de $x$. Se essas duas funções forem representadas por retas em um plano cartesiano (custo $y$ no eixo vertical e distância $x$ no eixo horizontal), assinale a alternativa que descreve corretamente essas retas.
\begin{enumerate}[label=(\alph*),nosep,leftmargin=*]
\item A reta que representa a empresa A passa pela origem (0,0) e tem coeficiente angular 2,5, sendo uma função proporcional; a reta da empresa B corta o eixo y no ponto (0,3) e tem coeficiente angular 2, não sendo proporcional.
\item A reta que representa a empresa B passa pela origem (0,0), pois seu coeficiente angular é 2; a reta da empresa A corta o eixo y no ponto (0,3), pois seu coeficiente angular é 2,5.
\item Ambas as retas passam pela origem (0,0), já que representam custo por quilômetro rodado; a reta de A tem coeficiente angular 2,5 e a de B tem coeficiente angular 2.
\item A reta da empresa A corta o eixo y no ponto (0, 2,5); a reta da empresa B corta o eixo y no ponto (0,2) e tem coeficiente angular 3.
\end{enumerate}

\textbf{Gabarito proposto:} A reta que representa a empresa A passa pela origem $(0,0)$ e tem coeficiente angular 2,5 (função proporcional); a reta da empresa B corta o eixo y no ponto $(0,3)$ e tem coeficiente angular 2 (função afim, não proporcional).

\textbf{Habilidade declarada:} EM13MAT401 --- Converter representações algébricas de funções polinomiais de 1º grau para representações geométricas no plano cartesiano, distinguindo os casos nos quais o comportamento é proporcional, recorrendo ou não a softwares ou aplicativos de álgebra e geometria dinâmica.
\end{samepage}

\noindent\rule{\linewidth}{0.4pt}
\textbf{Tem erro matemático?}\quad
\mbox{$\square$~não} \quad \mbox{$\square$~sim, aqui:} \underline{\hspace{5cm}}
\quad \mbox{$\square$~não sei dizer}

\textbf{Que nível de dificuldade esta questão tem para os seus alunos?}\quad
\mbox{$\square$~fácil} \quad \mbox{$\square$~média} \quad
\mbox{$\square$~difícil} \quad \mbox{$\square$~fora do alcance da turma}

\textbf{Você usaria esta questão?}\quad
\mbox{$\square$~aceita} \quad \mbox{$\square$~aceita com ajuste} \quad
\mbox{$\square$~recusada}

\textbf{Por quê?} \emph{(obrigatório quando não for ``aceita'')}

\underline{\hspace{\linewidth}}

\underline{\hspace{\linewidth}}
\vspace{0.6em}

\begin{samepage}
\subsection*{Q072}
Marina aplicou R\$\,8.000,00 em um investimento que rende juros compostos a uma taxa fixa de 1,5\% ao mês sobre o saldo do mês anterior. Ela decidiu que só fará o resgate no primeiro mês em que o montante acumulado for igual ou superior ao dobro do capital investido. Usando as aproximações $\log 2 \approx 0,301$ e $\log 1,015 \approx 0,00647$, determine o número mínimo de meses completos necessários para que Marina possa efetuar esse resgate.
\begin{enumerate}[label=(\alph*),nosep,leftmargin=*]
\item 47 meses
\item 46 meses
\item 67 meses
\item 48 meses
\end{enumerate}

\textbf{Gabarito proposto:} 47 meses

\textbf{Habilidade declarada:} EM13MAT303 --- Resolver e elaborar problemas envolvendo porcentagens em diversos contextos e sobre juros compostos, destacando o crescimento exponencial.
\end{samepage}

\noindent\rule{\linewidth}{0.4pt}
\textbf{Tem erro matemático?}\quad
\mbox{$\square$~não} \quad \mbox{$\square$~sim, aqui:} \underline{\hspace{5cm}}
\quad \mbox{$\square$~não sei dizer}

\textbf{Que nível de dificuldade esta questão tem para os seus alunos?}\quad
\mbox{$\square$~fácil} \quad \mbox{$\square$~média} \quad
\mbox{$\square$~difícil} \quad \mbox{$\square$~fora do alcance da turma}

\textbf{Você usaria esta questão?}\quad
\mbox{$\square$~aceita} \quad \mbox{$\square$~aceita com ajuste} \quad
\mbox{$\square$~recusada}

\textbf{Por quê?} \emph{(obrigatório quando não for ``aceita'')}

\underline{\hspace{\linewidth}}

\underline{\hspace{\linewidth}}
\vspace{0.6em}

\begin{samepage}
\subsection*{Q048}
Considere a progressão geométrica $(a_n)$, com $n \in \mathbb{N}^{*}$ (ou seja, $n = 1, 2, 3, \dots$), definida por $a_1 = 4$ e razão $q = 2$.

a) Determine a lei de uma função exponencial $f:\mathbb{R}\to\mathbb{R}$ tal que $f(n) = a_n$ para todo $n \in \mathbb{N}^{*}$, explicitando como o primeiro termo e a razão da PG se tornam os parâmetros dessa função.

b) Compare os domínios de $(a_n)$ e de $f$, explicando por que a progressão pode ser vista como a restrição da função $f$ a um subconjunto discreto de $\mathbb{R}$.

c) Determine o valor real de $x$ para o qual $f(x) = 24$. Em seguida, usando o domínio da PG, decida (com justificativa) se existe algum termo $a_n$ da progressão igual a 24.

\textbf{Gabarito proposto:} a) $f(x) = 4\cdot 2^{x-1}$ (equivalente a $2^{x+1}$), com $a_1=4$ e $q=2$ tornando-se, respectivamente, o valor inicial e a base da exponencial. b) $(a_n)$ tem domínio discreto $\mathbb{N}^{*}$ e é a restrição de $f$, que tem domínio contínuo $\mathbb{R}$, ao conjunto $\mathbb{N}^{*}$. c) $x=\log_2 12\approx 3{,}585$, que não é natural; logo não existe termo da PG igual a 24 (situa-se entre $a_3=16$ e $a_4=32$).

\textbf{Habilidade declarada:} EM13MAT508 --- Identificar e associar sequências numéricas (PG) a funções exponenciais de domínios discretos para análise de propriedades, incluindo dedução de algumas fórmulas e resolução de problemas.
\end{samepage}

\noindent\rule{\linewidth}{0.4pt}
\textbf{Tem erro matemático?}\quad
\mbox{$\square$~não} \quad \mbox{$\square$~sim, aqui:} \underline{\hspace{5cm}}
\quad \mbox{$\square$~não sei dizer}

\textbf{Que nível de dificuldade esta questão tem para os seus alunos?}\quad
\mbox{$\square$~fácil} \quad \mbox{$\square$~média} \quad
\mbox{$\square$~difícil} \quad \mbox{$\square$~fora do alcance da turma}

\textbf{Você usaria esta questão?}\quad
\mbox{$\square$~aceita} \quad \mbox{$\square$~aceita com ajuste} \quad
\mbox{$\square$~recusada}

\textbf{Por quê?} \emph{(obrigatório quando não for ``aceita'')}

\underline{\hspace{\linewidth}}

\underline{\hspace{\linewidth}}
\vspace{0.6em}

\begin{samepage}
\subsection*{Q090}
Uma progressão geométrica $(a_n)$, definida apenas para $n \in \mathbb{N}^*$ (isto é, $n = 1, 2, 3, \dots$), tem $a_1 = 4$ e $a_4 = 32$. Sabe-se que essa PG é a restrição ao conjunto $\{1,2,3,\dots\}$ de uma função exponencial $f:\mathbb{Z}\to\mathbb{R}$, dada por $f(n) = A\cdot B^n$, tal que $f(n) = a_n$ para todo $n \geq 1$. Determine o valor de $f(0)$, isto é, o valor que a função exponencial assume fora do domínio original da progressão.
\begin{enumerate}[label=(\alph*),nosep,leftmargin=*]
\item $f(0) = 4$
\item $f(0) = 2$
\item $f(0) = 8$
\item $f(0) = \dfrac{1}{2}$
\end{enumerate}

\textbf{Gabarito proposto:} f(0) = 2

\textbf{Habilidade declarada:} EM13MAT508 --- Identificar e associar sequências numéricas (PG) a funções exponenciais de domínios discretos para análise de propriedades, incluindo dedução de algumas fórmulas e resolução de problemas.
\end{samepage}

\noindent\rule{\linewidth}{0.4pt}
\textbf{Tem erro matemático?}\quad
\mbox{$\square$~não} \quad \mbox{$\square$~sim, aqui:} \underline{\hspace{5cm}}
\quad \mbox{$\square$~não sei dizer}

\textbf{Que nível de dificuldade esta questão tem para os seus alunos?}\quad
\mbox{$\square$~fácil} \quad \mbox{$\square$~média} \quad
\mbox{$\square$~difícil} \quad \mbox{$\square$~fora do alcance da turma}

\textbf{Você usaria esta questão?}\quad
\mbox{$\square$~aceita} \quad \mbox{$\square$~aceita com ajuste} \quad
\mbox{$\square$~recusada}

\textbf{Por quê?} \emph{(obrigatório quando não for ``aceita'')}

\underline{\hspace{\linewidth}}

\underline{\hspace{\linewidth}}
\vspace{0.6em}

\begin{samepage}
\subsection*{Q015}
Uma colônia de bactérias em um experimento de laboratório tem seu número de indivíduos, a partir do início da observação, dado por $N(t) = 100 \cdot 2^{t}$, em que $t$ é o tempo em horas ($t \geq 0$) e $N(t)$ é o número de bactérias. Um pesquisador quer saber, ao contrário, quanto tempo é necessário para que a colônia atinja um determinado número $N$ de bactérias, e por isso utiliza a função inversa de $N(t)$, que ele chama de $t(N)$.

a) Determine o domínio e a imagem de $N(t)$, considerando o contexto do experimento.

b) Encontre a expressão de $t(N)$ (a função inversa de $N(t)$) e determine seu domínio e sua imagem.

c) Compare o crescimento das duas funções: as duas são crescentes, mas uma cresce 'cada vez mais rápido' e a outra cresce 'cada vez mais devagar' à medida que a variável aumenta. Identifique qual é qual e justifique observando como varia o incremento de cada função.

d) Explique por que o domínio de $N(t)$ coincide com a imagem de $t(N)$, e a imagem de $N(t)$ coincide com o domínio de $t(N)$, relacionando esse fato com a posição dos gráficos das duas funções no plano cartesiano.

\textbf{Gabarito proposto:} a) $D_N=[0,+\infty)$, $Im_N=[100,+\infty)$. b) $t(N)=\log_2(N/100)$, com $D_t=[100,+\infty)$, $Im_t=[0,+\infty)$. c) $N(t)$ cresce cada vez mais rápido (convexa); $t(N)$ cresce cada vez mais devagar (côncava). d) Por serem funções inversas, seus gráficos são simétricos em relação à reta $y=x$, o que troca domínio e imagem entre as duas funções.

\textbf{Habilidade declarada:} EM13MAT403 --- Comparar e analisar as representações, em plano cartesiano, das funções exponencial e logarítmica para identificar as características fundamentais (domínio, imagem, crescimento) de cada uma, com ou sem apoio de tecnologias digitais, estabelecendo relações entre elas.
\end{samepage}

\noindent\rule{\linewidth}{0.4pt}
\textbf{Tem erro matemático?}\quad
\mbox{$\square$~não} \quad \mbox{$\square$~sim, aqui:} \underline{\hspace{5cm}}
\quad \mbox{$\square$~não sei dizer}

\textbf{Que nível de dificuldade esta questão tem para os seus alunos?}\quad
\mbox{$\square$~fácil} \quad \mbox{$\square$~média} \quad
\mbox{$\square$~difícil} \quad \mbox{$\square$~fora do alcance da turma}

\textbf{Você usaria esta questão?}\quad
\mbox{$\square$~aceita} \quad \mbox{$\square$~aceita com ajuste} \quad
\mbox{$\square$~recusada}

\textbf{Por quê?} \emph{(obrigatório quando não for ``aceita'')}

\underline{\hspace{\linewidth}}

\underline{\hspace{\linewidth}}
\vspace{0.6em}

\begin{samepage}
\subsection*{Q001}
Uma empresa de transporte por aplicativo cobra, em cada corrida, uma taxa fixa somada a um valor proporcional à distância percorrida. Um cliente fez uma corrida de 4 km e pagou R\$\,13,00. Em outro dia, o mesmo cliente percorreu 10 km e pagou R\$\,25,00. Sabe-se que o preço da corrida varia linearmente com a distância percorrida.

a) Determine a expressão algébrica que fornece o preço C, em reais, de uma corrida em função da distância d, em quilômetros.

b) Usando essa expressão, calcule quanto custará uma corrida de 15 km.

\textbf{Gabarito proposto:} C(d) = 2d + 5; para d = 15 km, o custo é R\$\,35,00.

\textbf{Habilidade declarada:} EM13MAT302 --- Resolver e elaborar problemas cujos modelos são as funções polinomiais de 1º e 2º graus, em contextos diversos, incluindo ou não tecnologias digitais.
\end{samepage}

\noindent\rule{\linewidth}{0.4pt}
\textbf{Tem erro matemático?}\quad
\mbox{$\square$~não} \quad \mbox{$\square$~sim, aqui:} \underline{\hspace{5cm}}
\quad \mbox{$\square$~não sei dizer}

\textbf{Que nível de dificuldade esta questão tem para os seus alunos?}\quad
\mbox{$\square$~fácil} \quad \mbox{$\square$~média} \quad
\mbox{$\square$~difícil} \quad \mbox{$\square$~fora do alcance da turma}

\textbf{Você usaria esta questão?}\quad
\mbox{$\square$~aceita} \quad \mbox{$\square$~aceita com ajuste} \quad
\mbox{$\square$~recusada}

\textbf{Por quê?} \emph{(obrigatório quando não for ``aceita'')}

\underline{\hspace{\linewidth}}

\underline{\hspace{\linewidth}}
\vspace{0.6em}

\begin{samepage}
\subsection*{Q064}
Uma grandeza positiva $C$ sofre, em cada uma de $n$ etapas sucessivas, um acréscimo percentual constante de taxa $i$ (com $0<i<1$): ao final de cada etapa, o valor da grandeza fica multiplicado por $(1+i)$ em relação ao valor obtido na etapa anterior. Esse é exatamente o modelo de capitalização por juros compostos.

\textbf{a)} Deduza, a partir da regra acima, a expressão geral $C_n$ do valor da grandeza após $n$ etapas, em função de $C$, $i$ e $n$. Em seguida, explique por que essa dependência de $C_n$ em relação a $n$ caracteriza um crescimento \textbf{exponencial}, e não um crescimento linear (do tipo somar, a cada etapa, um valor fixo igual a $iC$).

\textbf{b)} Um colega afirma que, se em duas etapas sucessivas a grandeza recebe acréscimos percentuais de taxas $i_1$ e $i_2$ (não necessariamente iguais), então a taxa percentual única $I$ equivalente às duas etapas combinadas é simplesmente $I = i_1+i_2$. Usando a expressão obtida no item (a), construa a fórmula correta de $I$ em função de $i_1$ e $i_2$, e mostre algebricamente que, sempre que $i_1>0$ e $i_2>0$, tem-se $I > i_1+i_2$.

\textbf{c)} Aplicando o modelo do item (a) com uma taxa constante de $i=20\%$ por etapa, determine o menor número inteiro de etapas $n$ necessário para que a grandeza mais que dobre de valor em relação ao valor inicial $C$.

\textbf{Gabarito proposto:} (a) $C_n = C(1+i)^n$, exponencial pois $n$ está no expoente e a razão entre termos consecutivos é constante e igual a $(1+i)$; (b) $I = i_1+i_2+i_1i_2$, que excede $i_1+i_2$ em $i_1i_2>0$; (c) $n=4$.

\textbf{Habilidade declarada:} EM13MAT303 --- Resolver e elaborar problemas envolvendo porcentagens em diversos contextos e sobre juros compostos, destacando o crescimento exponencial.
\end{samepage}

\noindent\rule{\linewidth}{0.4pt}
\textbf{Tem erro matemático?}\quad
\mbox{$\square$~não} \quad \mbox{$\square$~sim, aqui:} \underline{\hspace{5cm}}
\quad \mbox{$\square$~não sei dizer}

\textbf{Que nível de dificuldade esta questão tem para os seus alunos?}\quad
\mbox{$\square$~fácil} \quad \mbox{$\square$~média} \quad
\mbox{$\square$~difícil} \quad \mbox{$\square$~fora do alcance da turma}

\textbf{Você usaria esta questão?}\quad
\mbox{$\square$~aceita} \quad \mbox{$\square$~aceita com ajuste} \quad
\mbox{$\square$~recusada}

\textbf{Por quê?} \emph{(obrigatório quando não for ``aceita'')}

\underline{\hspace{\linewidth}}

\underline{\hspace{\linewidth}}
\vspace{0.6em}

\begin{samepage}
\subsection*{Q013}
Um professor propõe que os alunos investiguem, usando um software gráfico (ou construindo os gráficos manualmente em papel milimetrado), as funções $f(x) = 2^x$ e $g(x) = \log_2 x$.

a) Determine o domínio e o conjunto-imagem de $f$ e de $g$. Compare esses dois conjuntos (domínio de $f$ com imagem de $g$, e domínio de $g$ com imagem de $f$) e explique, matematicamente, por que essa relação ocorre.

b) Calcule $f(3)$ e $g(8)$; calcule também $f(0)$ e $g(1)$. Observando os pares de valores obtidos, descreva que tipo de simetria os gráficos de $f$ e $g$ apresentam no plano cartesiano.

c) Classifique $f$ e $g$ quanto ao crescimento (crescente ou decrescente) em todo o domínio de cada uma. Em seguida, utilizando os valores calculados no item (b), explique por que o fato de $f$ e $g$ serem ambas crescentes é compatível com a simetria identificada — ou seja, justifique por que essa simetria não transforma o crescimento de uma função em decrescimento na outra.

\textbf{Gabarito proposto:} a) $D(f)=\mathbb{R}=Im(g)$ e $D(g)=(0,+\infty)=Im(f)$, pois $g$ é a inversa de $f$. b) $f(3)=8, g(8)=3, f(0)=1, g(1)=0$: os gráficos de $f$ e $g$ são simétricos em relação à reta $y=x$. c) $f$ e $g$ são ambas crescentes; a simetria por $y=x$ preserva a ordem dos valores, por isso não inverte o crescimento.

\textbf{Habilidade declarada:} EM13MAT403 --- Comparar e analisar as representações, em plano cartesiano, das funções exponencial e logarítmica para identificar as características fundamentais (domínio, imagem, crescimento) de cada uma, com ou sem apoio de tecnologias digitais, estabelecendo relações entre elas.
\end{samepage}

\noindent\rule{\linewidth}{0.4pt}
\textbf{Tem erro matemático?}\quad
\mbox{$\square$~não} \quad \mbox{$\square$~sim, aqui:} \underline{\hspace{5cm}}
\quad \mbox{$\square$~não sei dizer}

\textbf{Que nível de dificuldade esta questão tem para os seus alunos?}\quad
\mbox{$\square$~fácil} \quad \mbox{$\square$~média} \quad
\mbox{$\square$~difícil} \quad \mbox{$\square$~fora do alcance da turma}

\textbf{Você usaria esta questão?}\quad
\mbox{$\square$~aceita} \quad \mbox{$\square$~aceita com ajuste} \quad
\mbox{$\square$~recusada}

\textbf{Por quê?} \emph{(obrigatório quando não for ``aceita'')}

\underline{\hspace{\linewidth}}

\underline{\hspace{\linewidth}}
\vspace{0.6em}

\begin{samepage}
\subsection*{Q074}
Uma empresa de transporte urbano inaugurou uma nova linha de ônibus. No primeiro mês de operação (mês 1), a linha transportou 340 passageiros. Devido à divulgação e à adesão gradual dos moradores, o número de passageiros transportados em cada mês superou o do mês anterior em exatamente 45 passageiros, formando uma progressão aritmética que se manteve mês após mês.

A diretoria da empresa sabe que essa evolução mensal pode ser descrita por uma função afim $f$, definida apenas para os números naturais $n \geq 1$ (o mês de operação), de modo que $f(n)$ forneça o número de passageiros transportados no mês $n$, coincidindo com os termos da progressão aritmética.

Com base nessa função afim, determine a partir de qual mês o número de passageiros transportados ultrapassará, pela primeira vez, 10.000 passageiros mensais.
\begin{enumerate}[label=(\alph*),nosep,leftmargin=*]
\item Mês 216
\item Mês 215
\item Mês 217
\item Mês 16
\end{enumerate}

\textbf{Gabarito proposto:} Mês 216 (alternativa A)

\textbf{Habilidade declarada:} EM13MAT507 --- Identificar e associar sequências numéricas (PA) a funções afins de domínios discretos para análise de propriedades, incluindo dedução de algumas fórmulas e resolução de problemas.
\end{samepage}

\noindent\rule{\linewidth}{0.4pt}
\textbf{Tem erro matemático?}\quad
\mbox{$\square$~não} \quad \mbox{$\square$~sim, aqui:} \underline{\hspace{5cm}}
\quad \mbox{$\square$~não sei dizer}

\textbf{Que nível de dificuldade esta questão tem para os seus alunos?}\quad
\mbox{$\square$~fácil} \quad \mbox{$\square$~média} \quad
\mbox{$\square$~difícil} \quad \mbox{$\square$~fora do alcance da turma}

\textbf{Você usaria esta questão?}\quad
\mbox{$\square$~aceita} \quad \mbox{$\square$~aceita com ajuste} \quad
\mbox{$\square$~recusada}

\textbf{Por quê?} \emph{(obrigatório quando não for ``aceita'')}

\underline{\hspace{\linewidth}}

\underline{\hspace{\linewidth}}
\vspace{0.6em}

\begin{samepage}
\subsection*{Q005}
Uma pessoa comparou dois aplicativos de mototáxi usando cada um deles em meses diferentes, sempre para o mesmo tipo de deslocamento no trabalho. Em cada aplicativo, a fatura mensal é composta por uma tarifa fixa mais um valor constante por quilômetro total rodado no mês (isto é, o custo mensal é uma função afim da quilometragem mensal $x$, em km, com $0 \le x \le 150$). Os registros de fatura foram:

\begin{itemize}[nosep,leftmargin=*]
\item MotoVá: em um mês, rodou 40 km e pagou R\$\,62,00; em outro mês, rodou 100 km e pagou R\$\,122,00.
\item UrbanMoto: em um mês, rodou 40 km e pagou R\$\,70,00; em outro mês, rodou 90 km e pagou R\$\,100,00.
\end{itemize}

a) Determine a lei $c_A(x)$ que dá o custo mensal do MotoVá e a lei $c_B(x)$ que dá o custo mensal do UrbanMoto, em função da quilometragem mensal $x$.

b) Determine a quilometragem mensal para a qual os dois aplicativos cobram exatamente o mesmo valor, e qual é esse valor. Indique qual aplicativo é mais barato para quilometragens menores que esse valor e qual é mais barato para quilometragens maiores.

c) Um terceiro aplicativo, o RotaFácil, acaba de entrar no mercado, cobrando uma tarifa fixa de R\$\,50,00 mais R\$\,0,50 por quilômetro rodado no mês. Elabore uma única função $C(x)$, definida por partes para $0 \le x \le 150$, que forneça sempre o menor custo mensal possível entre os três aplicativos, para qualquer quilometragem mensal $x$ nesse intervalo (explicite os intervalos de $x$ de cada sentença). Em seguida, justifique com os cálculos por que o UrbanMoto não aparece em nenhuma sentença da função que você construiu, ou seja, mostre que ele nunca é a opção mais barata para nenhum valor de $x$ em $[0,150]$.

\textbf{Gabarito proposto:} c\_A(x) = x + 22; c\_B(x) = 0,6x + 46; equilíbrio A–B em x = 60 km (custo R\$\,82,00), com MotoVá mais barato para x<60 e UrbanMoto mais barato para x>60. Com o RotaFácil, c\_C(x) = 0,5x + 50, e a função de menor custo é C(x) = x+22 para 0$\le$x$\le$56 e C(x) = 0,5x+50 para 56<x$\le$150; o UrbanMoto nunca é a opção mais barata em [0,150].

\textbf{Habilidade declarada:} EM13MAT302 --- Resolver e elaborar problemas cujos modelos são as funções polinomiais de 1º e 2º graus, em contextos diversos, incluindo ou não tecnologias digitais.
\end{samepage}

\noindent\rule{\linewidth}{0.4pt}
\textbf{Tem erro matemático?}\quad
\mbox{$\square$~não} \quad \mbox{$\square$~sim, aqui:} \underline{\hspace{5cm}}
\quad \mbox{$\square$~não sei dizer}

\textbf{Que nível de dificuldade esta questão tem para os seus alunos?}\quad
\mbox{$\square$~fácil} \quad \mbox{$\square$~média} \quad
\mbox{$\square$~difícil} \quad \mbox{$\square$~fora do alcance da turma}

\textbf{Você usaria esta questão?}\quad
\mbox{$\square$~aceita} \quad \mbox{$\square$~aceita com ajuste} \quad
\mbox{$\square$~recusada}

\textbf{Por quê?} \emph{(obrigatório quando não for ``aceita'')}

\underline{\hspace{\linewidth}}

\underline{\hspace{\linewidth}}
\vspace{0.6em}

\begin{samepage}
\subsection*{Q016}
Observe a tabela abaixo, que relaciona valores de uma grandeza $x$ a valores de uma grandeza $y$, obtidos a partir de um mesmo padrão:

\begin{center}
\begin{tabular}{|p{0.172\linewidth}|p{0.172\linewidth}|p{0.172\linewidth}|p{0.172\linewidth}|p{0.172\linewidth}|}
\hline
\textbf{$x$} & \textbf{1} & \textbf{2} & \textbf{3} & \textbf{4} \\
\hline
$y$ & 5 & 8 & 11 & 14 \\
\hline
\end{tabular}
\end{center}

a) Analise como $y$ varia à medida que $x$ aumenta uma unidade e identifique o padrão presente na tabela.

b) Escreva a lei algébrica $y = f(x)$ que generaliza esse padrão para qualquer valor real de $x$.

c) Classifique o tipo de função obtida, justificando sua resposta a partir da lei encontrada.

\textbf{Gabarito proposto:} A relação é uma função afim dada por $f(x) = 3x + 2$, pois a variação de $y$ é constante (igual a 3) para cada acréscimo unitário de $x$.

\textbf{Habilidade declarada:} EM13MAT501 --- Investigar relações entre números expressos em tabelas para representá-los no plano cartesiano, identificando padrões e criando conjecturas para generalizar e expressar algebricamente essa generalização, reconhecendo quando essa representação é de função polinomial de 1º grau.
\end{samepage}

\noindent\rule{\linewidth}{0.4pt}
\textbf{Tem erro matemático?}\quad
\mbox{$\square$~não} \quad \mbox{$\square$~sim, aqui:} \underline{\hspace{5cm}}
\quad \mbox{$\square$~não sei dizer}

\textbf{Que nível de dificuldade esta questão tem para os seus alunos?}\quad
\mbox{$\square$~fácil} \quad \mbox{$\square$~média} \quad
\mbox{$\square$~difícil} \quad \mbox{$\square$~fora do alcance da turma}

\textbf{Você usaria esta questão?}\quad
\mbox{$\square$~aceita} \quad \mbox{$\square$~aceita com ajuste} \quad
\mbox{$\square$~recusada}

\textbf{Por quê?} \emph{(obrigatório quando não for ``aceita'')}

\underline{\hspace{\linewidth}}

\underline{\hspace{\linewidth}}
\vspace{0.6em}

\begin{samepage}
\subsection*{Q055}
Em um teste de frenagem, a distância de frenagem $d$ (em metros) de um carrinho de brinquedo, em função de sua velocidade $v$ (em m/s, com $v \geq 0$), foi modelada por $d(v) = 0{,}5v^2$.

Em outro experimento, a altura $h$ (em metros) de uma bolinha lançada verticalmente para cima a partir de uma plataforma, em função do tempo $t$ (em segundos, com $t \geq 0$), foi modelada por $h(t) = -5t^2 + 20t + 1$.

Um estudante deseja esboçar, sem construir tabela de valores, o gráfico no plano cartesiano correspondente à situação em que a grandeza dependente é diretamente proporcional ao quadrado da grandeza independente.

Assinale a alternativa que descreve corretamente esse gráfico.
\begin{enumerate}[label=(\alph*),nosep,leftmargin=*]
\item Parábola com concavidade voltada para cima, vértice na origem $(0,0)$, simétrica em relação ao eixo vertical e situada no primeiro quadrante para $v \geq 0$, pois $d(v)=0{,}5v^2$ é diretamente proporcional ao quadrado de $v$.
\item Parábola com concavidade voltada para baixo, vértice no ponto $(2,21)$, interceptando o eixo vertical em $(0,1)$.
\item Reta crescente passando pela origem, com coeficiente angular $0{,}5$.
\item Parábola com concavidade voltada para cima, vértice no ponto $(0; 0{,}5)$, deslocada verticalmente para cima em relação à origem.
\end{enumerate}

\textbf{Gabarito proposto:} A

\textbf{Habilidade declarada:} EM13MAT402 --- Converter representações algébricas de funções polinomiais de 2º grau para representações geométricas no plano cartesiano, distinguindo os casos nos quais uma variável for diretamente proporcional ao quadrado da outra, recorrendo ou não a softwares ou aplicativos de álgebra e geometria dinâmica.
\end{samepage}

\noindent\rule{\linewidth}{0.4pt}
\textbf{Tem erro matemático?}\quad
\mbox{$\square$~não} \quad \mbox{$\square$~sim, aqui:} \underline{\hspace{5cm}}
\quad \mbox{$\square$~não sei dizer}

\textbf{Que nível de dificuldade esta questão tem para os seus alunos?}\quad
\mbox{$\square$~fácil} \quad \mbox{$\square$~média} \quad
\mbox{$\square$~difícil} \quad \mbox{$\square$~fora do alcance da turma}

\textbf{Você usaria esta questão?}\quad
\mbox{$\square$~aceita} \quad \mbox{$\square$~aceita com ajuste} \quad
\mbox{$\square$~recusada}

\textbf{Por quê?} \emph{(obrigatório quando não for ``aceita'')}

\underline{\hspace{\linewidth}}

\underline{\hspace{\linewidth}}
\vspace{0.6em}

\begin{samepage}
\subsection*{Q007}
Uma oficina mecânica cobra por seus serviços de revisão um valor fixo de mão de obra, mais uma taxa que depende do número de horas de trabalho. Uma revisão que durou 2 horas custou R\$\,150,00, e outra revisão, no mesmo padrão de cobrança, durou 5 horas e custou R\$\,300,00. Sabendo que o custo total varia linearmente com o número de horas trabalhadas, qual seria o custo de uma revisão que durasse 8 horas?
\begin{enumerate}[label=(\alph*),nosep,leftmargin=*]
\item R\$\,450,00
\item R\$\,400,00
\item R\$\,600,00
\item R\$\,480,00
\end{enumerate}

\textbf{Gabarito proposto:} R\$\,450,00

\textbf{Habilidade declarada:} EM13MAT302 --- Resolver e elaborar problemas cujos modelos são as funções polinomiais de 1º e 2º graus, em contextos diversos, incluindo ou não tecnologias digitais.
\end{samepage}

\noindent\rule{\linewidth}{0.4pt}
\textbf{Tem erro matemático?}\quad
\mbox{$\square$~não} \quad \mbox{$\square$~sim, aqui:} \underline{\hspace{5cm}}
\quad \mbox{$\square$~não sei dizer}

\textbf{Que nível de dificuldade esta questão tem para os seus alunos?}\quad
\mbox{$\square$~fácil} \quad \mbox{$\square$~média} \quad
\mbox{$\square$~difícil} \quad \mbox{$\square$~fora do alcance da turma}

\textbf{Você usaria esta questão?}\quad
\mbox{$\square$~aceita} \quad \mbox{$\square$~aceita com ajuste} \quad
\mbox{$\square$~recusada}

\textbf{Por quê?} \emph{(obrigatório quando não for ``aceita'')}

\underline{\hspace{\linewidth}}

\underline{\hspace{\linewidth}}
\vspace{0.6em}

\newpage
\subsection*{Para terminar}

\textbf{1. Você usaria um sistema assim no seu planejamento? Por quê?}

\underline{\hspace{\linewidth}}

\underline{\hspace{\linewidth}}

\bigskip
\textbf{2. O que faltou nas questões que você viu?}

\underline{\hspace{\linewidth}}

\underline{\hspace{\linewidth}}

\bigskip
\textbf{3. O que atrapalhou --- na questão, no enunciado, no formato deste
material?}

\underline{\hspace{\linewidth}}

\underline{\hspace{\linewidth}}

\bigskip
Obrigado. Se quiser saber quais itens tinham erro e qual era, é só pedir.
\end{document}